% Options for packages loaded elsewhere
\PassOptionsToPackage{unicode}{hyperref}
\PassOptionsToPackage{hyphens}{url}
%
\documentclass[
]{article}
\usepackage{amsmath,amssymb}
\usepackage{iftex}
\ifPDFTeX
  \usepackage[T1]{fontenc}
  \usepackage[utf8]{inputenc}
  \usepackage{textcomp} % provide euro and other symbols
\else % if luatex or xetex
  \usepackage{unicode-math} % this also loads fontspec
  \defaultfontfeatures{Scale=MatchLowercase}
  \defaultfontfeatures[\rmfamily]{Ligatures=TeX,Scale=1}
\fi
\usepackage{lmodern}
\ifPDFTeX\else
  % xetex/luatex font selection
\fi
% Use upquote if available, for straight quotes in verbatim environments
\IfFileExists{upquote.sty}{\usepackage{upquote}}{}
\IfFileExists{microtype.sty}{% use microtype if available
  \usepackage[]{microtype}
  \UseMicrotypeSet[protrusion]{basicmath} % disable protrusion for tt fonts
}{}
\makeatletter
\@ifundefined{KOMAClassName}{% if non-KOMA class
  \IfFileExists{parskip.sty}{%
    \usepackage{parskip}
  }{% else
    \setlength{\parindent}{0pt}
    \setlength{\parskip}{6pt plus 2pt minus 1pt}}
}{% if KOMA class
  \KOMAoptions{parskip=half}}
\makeatother
\usepackage{xcolor}
\usepackage[margin=1in]{geometry}
\usepackage{graphicx}
\makeatletter
\def\maxwidth{\ifdim\Gin@nat@width>\linewidth\linewidth\else\Gin@nat@width\fi}
\def\maxheight{\ifdim\Gin@nat@height>\textheight\textheight\else\Gin@nat@height\fi}
\makeatother
% Scale images if necessary, so that they will not overflow the page
% margins by default, and it is still possible to overwrite the defaults
% using explicit options in \includegraphics[width, height, ...]{}
\setkeys{Gin}{width=\maxwidth,height=\maxheight,keepaspectratio}
% Set default figure placement to htbp
\makeatletter
\def\fps@figure{htbp}
\makeatother
\setlength{\emergencystretch}{3em} % prevent overfull lines
\providecommand{\tightlist}{%
  \setlength{\itemsep}{0pt}\setlength{\parskip}{0pt}}
\setcounter{secnumdepth}{5}
\usepackage{fancyhdr}
\pagestyle{fancy}
\fancyhead[L]{Hidden Markov Models for Stock Market Regime Detection}
\fancyfoot[CO,CE]{}
\fancyfoot[LE,RO]{\thepage}
\usepackage{setspace}\doublespacing
\ifLuaTeX
  \usepackage{selnolig}  % disable illegal ligatures
\fi
\usepackage{bookmark}
\IfFileExists{xurl.sty}{\usepackage{xurl}}{} % add URL line breaks if available
\urlstyle{same}
\hypersetup{
  pdftitle={Hidden Markov Models for Stock Market Regime Detection},
  pdfauthor={Jayant Saini, Udayan Sablok, Pranshu Sharma},
  hidelinks,
  pdfcreator={LaTeX via pandoc}}

\title{Hidden Markov Models for Stock Market Regime Detection}
\author{Jayant Saini, Udayan Sablok, Pranshu Sharma}
\date{2026-11-19}

\begin{document}
\maketitle

\newpage 
\tableofcontents 
\newpage

\newpage

\section{Introduction}\label{introduction}

Financial markets do not behave consistently over time, as periods of
growth, decline, stability, and high volatility can occur under
different market conditions. To describe these market conditions more
precisely, we define three market regimes: bull, bear and
consolidation/sideways. A bear market regime occurs when stock prices
decline by 20\% or more from a recent peak over a sustained period of
time. A bull market regime occurs when stock prices rise 20\% or more
from a recent low over a sustained period of time. A sideways market
regime occurs when stock prices fluctuate up and down within a narrow
range. An important distinction is that these market regimes are not
directly observable. We can observe stock prices and returns, but we
cannot directly observe whether the market is currently in a bull, bear,
or sideways regime.

Due to the ``hidden'' nature of the stock market regimes, Hidden Markov
Models provide a helpful framework for identifying and analysing these
unobservable regimes. Hidden Markov models can be used to identify the
likelihood of moving between market regimes and the signals given by the
market (stock-prices) that define each state. The primary objective of
this project is to use both simulated and real-world stock-market data
to fit a Hidden Markov Model, and interpret the characteristics and
transition probabilities of the regimes it identifies.

The rest of this project is as follows. Section 2 reviews the theory of
Markov chains and Hidden Markov Models. Section 3 presents the
simulation study and real stock-market data application (all code and
outputs will be in an Appendix). Section 4 summarizes the main findings
and limitations of the project, followed by References in section 5.

\section{Review of the Topic}\label{review-of-the-topic}

A Hidden Markov Model is a model where we observe a sequence of
emissions, in this case stock market returns; however, the sequence of
states the model went through to generate these emissions is unknown,
i.e., hidden. In our project, the hidden states are stock market
regimes, as we cannot directly observe them but must instead infer them
from stock market data. Hidden Markov Models are built on the theory of
Markov Chains, i.e.~the hidden states follow the Markov assumption. To
understand how HMMs can be applied to stock market regime detection, we
first review the fundamentals of Markov chains, then introduce Hidden
Markov Models and their applications to stock markets.

\subsection{Markov Chains}\label{markov-chains}

A markov chain models transition between states under the assumption of
the Markov property. \[
P(X_{t+1}=j\mid X_t=i).
\] The transition probabilties are modelled by the transition matrix (P)
which is defined as: \[
P_{ij}=P(X_{t+1}=j\mid X_t=i),
\qquad
\sum_j P_{ij}=1.
\] For the purpose of this project, we will be using a finite state,
time-homogeneous Markov chain. The states in the Markov Chain will
represent stock market regimes i.e.~bull, bear and sideways
(finite-state Markov chain). The transition matrix for our project will
define the probabilities of moving between these states over time.

\subsection{Hidden Markov Models}\label{hidden-markov-models}

Hidden Markov Models are an extension of the Markov Chain framework
where the underlying states are unobservable i.e.~hidden and can only be
inferred through observable data. Each of the underlying ``hidden''
states generates output according to a probability distribution which is
known as the emission distribution. The outputs generated by the states
are directly observable. In our project the hidden states we aim to
identify are as follows: bull, bear, and sideways, i.e.~the hidden
states represent market regimes, and the observations from each hidden
state we collect are stock market returns.

Hidden Markov Models have three main components:

\begin{enumerate}
\item \textbf{Transition Matrix (A):} Represents the probabilities of transitioning between hidden states.
\item \textbf{Emission Distribution (B):} Describes the probability distribution of stock market returns (observed values) within each hidden state.
\item \textbf{Initial Distribution ($\pi$):} Represents the probability of starting in a hidden regime.
\end{enumerate}

These three components together define a Hidden Markov Model: \[
\lambda = (A, B, \pi).
\]

Visually, the Hidden Markov Model looks as follows:

\begin{figure}

{\centering \includegraphics{Report_Template_files/figure-latex/hmm-diagram-1} 

}

\caption{Structure of a Hidden Markov Model}\label{fig:hmm-diagram}
\end{figure}

The horizontal arrows represent transitions between hidden states which
are based on the transition matrix A. The vertical arrows represent
observations generated by each state according to the emission
distribution B.

Since stock market returns are continuous observations, our project
models the emission distributions using a Gaussian distribution. In the
Gaussian distribution each hidden state i.e.~market regime has a mean
and variance. The use of the Gaussian emission distribution for hidden
states is defined mathematically as: \[
Y_t \mid X_t=i \sim \mathcal{N}(\mu_i,\sigma_i^2)
\]

\begin{itemize}
\item $Y_t$: observed stock market return at time $t$.
\item $X_t=i$: the hidden state $i$ the market is currently in.
\item $\mu_i$: mean of stock market returns in state $i$.
\item $\sigma_i^2$: variance of stock market returns in state $i$.
\end{itemize}

\subsection{Details of the topic}\label{details-of-the-topic}

This is a mathematical formua without formula number.
\[h(x) = \frac{f(x)}{S(x)} = -\frac{d \log[S(x])}{\rm dx}\]

This is how write a fomula and label it. \begin{align}
\label{eq2}
S(x) = \exp[-H(x)] = \exp[-\int_0^x h(u) {\rm d}u]
\end{align}

\section{Application}\label{application}

This is an R Markdown document. Markdown is a simple formatting syntax
for authoring HTML, PDF, and MS Word documents. For more details on
using R Markdown see \url{http://rmarkdown.rstudio.com}.

\begin{table}[!htp]
\small
\caption{This is a caption.}
\begin{center}
\begin{tabular}{|c|cc|}
\hline
& {Placebo}   &  {Treatment}       \\ \cline{1-3}  
Method& {Est} (SE)    &{Est} (SE) \\    %  95\% CI \hspace{0.5cm}  95\% CI 
\hline
\scriptsize{Method 1 } & ??? (???)     & ??? (???)      \\   %   (???, ???) \hspace{0.25cm} (???, ???)
\scriptsize{Method 2}  &  ??? (???)     &  ??? (???)     \\  %  (???,  ???)  \hspace{0.25cm}  (???, ???)
\hline
\end{tabular}
\label{tab: 2.3}
\end{center}
\end{table}

\section{Conclusion}\label{conclusion}

\section{References}\label{references}

Borst, David. ``Detecting Market Regimes: Hidden Markov Model.''
\emph{Medium}, 26 Aug.~2025,
\url{https://datadave1.medium.com/detecting-market-regimes-hidden-markov-model-2462e819c72e}.

Hayes, Adam. ``Sideways Markets Explained: Trading Tactics and
Benefits.'' \emph{Investopedia}, 21 May 2026,
\url{https://www.investopedia.com/terms/s/sidewaysmarket.asp}.

The MathWorks, Inc.~(n.d.). \emph{Hidden Markov models (HMM).} MATLAB \&
Simulink. Retrieved November 16, 2026, from
\url{https://es.mathworks.com/help/stats/hidden-markov-models-hmm.html}

QuestDB. ``Market Regime Detection Using Hidden Markov Models.''
\emph{QuestDB Glossary},
\url{https://questdb.com/glossary/market-regime-detection-using-hidden-markov-models/}.

Scotiabank Advice+ Team. ``Bull vs.~Bear Market: The Basics Explained.''
\emph{Scotiabank}, 27 Oct.~2026,
\url{https://www.scotiabank.com/ca/en/personal/advice-plus/features/posts.bear-vs-bull-market-the-basics.html}.

Stanford Institute for Human-Centered Artificial Intelligence. (n.d.).
\emph{What is a Markov chain?} Stanford HAI. Retrieved November 16,
2026, from
\url{https://hai.stanford.edu/ai-definitions/what-is-a-markov-chain}

\end{document}
